\documentclass[10pt,twoside]{article}
\usepackage[utf8]{inputenc}
\usepackage{amsmath}
\usepackage{amsfonts}
\usepackage{amssymb}
\usepackage[spanish,es-noshorthands]{babel}
\usepackage[T1]{fontenc}
\usepackage{lmodern}
\usepackage{graphicx,hyperref}
\usepackage{tikz,pgf}
\usepackage{marvosym}
\usepackage{multicol}
\usepackage{subfig}
\usepackage[papersize={5.5in,8.5in},total={4.75in,7.25in},centering]{geometry}
\usepackage{fancyhdr}
\pagestyle{fancy}
\fancyhead[LE]{IEDAB}
\fancyhead[LO]{PEI: ``Hacia una cultura del desarrollo sostenible''}
\fancyfoot[RE]{\url{www.autistici.org/mathgerman}}
\fancyfoot[RO]{\Email~iedabgerman@autistici.org}

\author{Germ\'an Avenda\~no Ram\'irez~\thanks{Lic. Mat. U.D., M.Sc. U.N.}}
\title{\begin{minipage}{.2\textwidth}
\includegraphics[height=1.75cm]{Images/logo-colegio.png}\end{minipage}
\begin{minipage}{.55\textwidth}
\begin{center}
Taller Combinatoria\\
Probabilidad $11^{\circ}$
\end{center}
\end{minipage}\hfill
\begin{minipage}{.2\textwidth}
\includegraphics[height=1.75cm]{Images/logo-sed.png} 
\end{minipage}}
\date{}
\begin{document}
\maketitle
\section*{Permutaciones}
El número de permutaciones de $n$ objetos de los cuales $n_{1}$ son iguales, $n_{2}$ son iguales, . . . es
\[\dfrac{n!}{n_{1}!n_{2}!\cdot \ldots} \qquad \text{donde }\; n=n_{1}+n_{2}+\ldots\]
\subsection*{Ejemplo:}
El número de permutaciones de las letras en la palabra statistics es:
\[\dfrac{10!}{3!3!1!2!1!}=50\,400\] porque hay 3 eses, 3 tes, 1 a, 2 íes y 1 c.
\subsection*{Ejercicios}
En los ejercicios siguientes cuente el número de formas en que puede hacerse cada procedimiento.
\begin{enumerate}
\item Alinear a tres personas para una fotografía.
\item Acomodar en línea 5 canicas de colores diferentes
\item Sentar a 10 personas en una banca en la que sólo hay 4 asientos disponibles
\item Sentar en hilera a 5 hombres y 4 mujeres de manera que las mujeres ocupen los lugares pares.
\item Los números de cuatro dígitos que se pueden formar con los 10 dígitos 0, 1, 2, 3, \ldots , 9, si: a) puede haber
repeticiones, b) no puede haber repeticiones y c) no puede haber repeticiones y el último dígito debe ser cero?
\item En un librero se van a acomodar cuatro libros diferentes de matemáticas, 6 libros diferentes de física y 2 libros diferentes de química. Cómo se acomodan si: a) los libros de cada materia tienen que estar juntos y b) sólo los libros de matemáticas tienen que estar juntos
\item Cinco canicas rojas, 2 canicas blancas y 3 azules están ordenadas en línea. Si las canicas de un mismo color no se distinguen unas de otras.
\item El número de formas de sentarse 7 personas a una mesa redonda si: a) las 7 se pueden sentar en cualquier lugar y b) 2 determinadas personas no pueden sentarse juntas
\section*{Combinaciones}
\item ¿De cuántas maneras pueden colocarse 10 objetos en dos grupos, uno de 4 y otro de 6 objetos?
\item ¿De cuántas maneras puede formarse de un grupo de 9 personas un comité de 5 personas?
\item Con 5 matemáticos y 7 físicos hay que formar un comité que conste de 2 matemáticos y 3 físicos. ¿De cuántas maneras se puede formar este comité si: a) puede incluirse a cualquiera de los matemáticos y a cualquiera de los físicos, b) hay uno de los físicos que tiene que formar parte del comité y c) hay dos de los matemáticos que no pueden formar parte del comité?
\item Una niña tiene 5 flores que son todas distintas. ¿Cuántos ramos puede formar?
\item Con 7 consonantes y 5 vocales ¿cuántas palabras con 4 consonantes distintas y 3 vocales distintas pueden formarse? No importa que las palabras no tengan significado.
\section*{Probabilidad y análisis combinatorio}
\item Una caja contiene 8 pelotas rojas, 3 blancas y 9 azules. Si se extraen 3 pelotas en forma aleatoria, determinar
la probabilidad de que: a) las 3 sean rojas, b) las 3 sean blancas, c) 2 sean rojas y 1 sea blanca, d ) por lo menos 1 sea blanca, e) se extraiga una de cada color y f ) se extraigan en el orden roja, blanca, azul.
\item De una baraja de 52 cartas bien barajadas se extraen 5 cartas. Encontrar la probabilidad de que: a) 4 sean ases; b) 4 sean ases y 1 sea rey; c) 3 sean dieces y 2 sean sotas; d ) sean 9, 10, sota, reina y rey en cualquier orden; e) 3 sean de un palo y 2 de otro palo, y f ) se obtenga por lo menos 1 as.
\item Determinar la probabilidad de tener 3 seises en cinco lanzamientos de un dado.
%\item En una fábrica se encuentra que en promedio 20\% de los tornillos producidos con una máquina están defectuosos. Si se toman aleatoriamente 10 tornillos producidos con esta máquina en un día, encontrar la probabilidad de que: a) exactamente 2 estén defectuosos, b) 2 o más estén defectuosos y c) más de 5 estén defectuosos.
\end{enumerate}

\end{document}
